\clearpage
\phantomsection
\addcontentsline{toc}{chapter}{Shodhsar (Abstract in Nepali)}
{\nepalifont
\begin{center}
{\fontsize{18}{22}\selectfont शोधसार}
\end{center}

\vspace{1.5\baselineskip}
\fontsize{14}{20}\selectfont

कृषि वैद्य {\normalfont B.Sc. CSIT} अन्तिम वर्षको परियोजनाका रूपमा विकास गरिएको
एकीकृत बाली स्वास्थ्य निदान तथा सल्लाह प्रणाली हो। यस परियोजनामा
{\normalfont YOLO} आधारित वस्तु पहिचान मोडेल, {\normalfont Qwen} आधारित
भिजन-ल्यान्ग्वेज मोडेल, {\normalfont Express/Mongoose} ब्याकएन्ड
{\normalfont API}, र {\normalfont React Native/Expo} मोबाइल अनुप्रयोगलाई एउटै
निदान कार्यप्रवाहमा जोडिएको छ। किसानले बालीको तस्वीर मोबाइल अनुप्रयोगमार्फत
पठाउँदा प्रणालीले ब्याकएन्ड र मोडेल सेवाको सहायताले संरचित निदान तथा सल्लाह
उपलब्ध गराउने उद्देश्य राखेको छ।

{\normalfont YOLO} भाग कृषि बोटबिरुवाका भागहरू पहिचान गर्न ३०,९८१ तस्वीर र
१,४७,४६७ {\normalfont annotation} भएको {\normalfont dataset} प्रयोग गरी तयार
गरिएको हो। {\normalfont VLM} भाग रोगसम्बन्धी व्याख्या उत्पादन गर्न
{\normalfont structured image-text data}, {\normalfont LoRA fine-tuning},
{\normalfont quantized training}, र {\normalfont prediction parsing} प्रयोग गरी
तयार गरिएको हो। {\normalfont Backend} ले {\normalfont authentication},
{\normalfont data model}, {\normalfont repository}, {\normalfont scan handling},
{\normalfont diagnosis orchestration}, {\normalfont validation middleware}, र
{\normalfont advisory response} सम्हाल्छ। {\normalfont Mobile app} ले
{\normalfont scan workflow}, {\normalfont crop record},
{\normalfont local SQLite persistence}, {\normalfont synchronization structure},
{\normalfont localization}, र {\normalfont farmer-facing result display} समेट्छ।

परियोजनाले {\normalfont model metrics}, {\normalfont validation outputs},
{\normalfont architecture diagrams}, {\normalfont source-backed tables}, र
{\normalfont documented tests} मार्फत कार्यान्वयन भएको
{\normalfont integrated prototype} प्रस्तुत गर्छ। मुख्य सीमितताहरू
{\normalfont VLM prediction reliability}, {\normalfont live end-to-end runtime evidence}
को अभाव, {\normalfont mobile device validation} को अपूर्णता, र
{\normalfont real-world advisory} प्रयोग अघि {\normalfont agricultural expert review}
आवश्यक हुनु हुन्।
}

\vspace{2\baselineskip}
\noindent\textit{\fontsize{10}{12}\selectfont Keywords:} {\fontsize{10}{12}\selectfont Krishi Vaidya, crop disease diagnosis, YOLO, vision-language model, mobile agriculture, backend API}
